% !TeX program = pdflatex
\documentclass[11pt,a4paper]{amsart}

\newcommand{\notelanguage}{finnish}
% Shared preamble for course notes and exercise sets.

\usepackage[T1]{fontenc}
\usepackage{lmodern}
\providecommand{\notelanguage}{english}
\usepackage[finnish,english]{babel}
\AtBeginDocument{\expandafter\selectlanguage\expandafter{\notelanguage}}

\newcommand{\notesfinnishlanguage}{finnish}
\ifx\notelanguage\notesfinnishlanguage
\newcommand{\notestheoremname}{Lause}
\newcommand{\noteslemmaname}{Lemma}
\newcommand{\notespropositionname}{Propositio}
\newcommand{\notescorollaryname}{Korollaari}
\newcommand{\notesdefinitionname}{Määritelmä}
\newcommand{\notesexamplename}{Esimerkki}
\newcommand{\notesproblemname}{Tehtävä}
\newcommand{\notesexercisename}{Tehtävä}
\newcommand{\notesremarkname}{Huomautus}
\newcommand{\notessolutionname}{Ratkaisu}
\else
\newcommand{\notestheoremname}{Theorem}
\newcommand{\noteslemmaname}{Lemma}
\newcommand{\notespropositionname}{Proposition}
\newcommand{\notescorollaryname}{Corollary}
\newcommand{\notesdefinitionname}{Definition}
\newcommand{\notesexamplename}{Example}
\newcommand{\notesproblemname}{Problem}
\newcommand{\notesexercisename}{Exercise}
\newcommand{\notesremarkname}{Remark}
\newcommand{\notessolutionname}{Solution}
\fi

\usepackage[a4paper,margin=30mm]{geometry}
\usepackage{microtype}
\usepackage{graphicx}
\usepackage{enumitem}

\newcommand{\exerciseheading}[2]{%
  \par\noindent
  {\large\bfseries #1\par}
  \vspace{0.1em}
  \noindent{\large #2\par}
  \vspace{1.5em}
}

\usepackage{amsmath}
\usepackage{amssymb}
\usepackage{amsthm}
\usepackage{mathtools}
\usepackage{aliascnt}

% Row operations beside the affected row, e.g. \rowop{+2R_1}.
\newcommand{\rowop}[1]{\mathrel{\big|}\,#1}

% Linear systems: the argument gives the number of alignment pairs.
\newenvironment{linearsystem}[1]
{\left\{
  \begin{alignedat}{#1}}
    {
  \end{alignedat}\right.}

  \usepackage[hidelinks,pdfusetitle]{hyperref}

  \numberwithin{equation}{section}
  \setlist{itemsep=0.25em,topsep=0.5em}

  \theoremstyle{plain}
  \newtheorem{theorem}{\notestheoremname}[section]
  \newaliascnt{lemma}{theorem}
  \newtheorem{lemma}[lemma]{\noteslemmaname}
  \aliascntresetthe{lemma}
  \newaliascnt{proposition}{theorem}
  \newtheorem{proposition}[proposition]{\notespropositionname}
  \aliascntresetthe{proposition}
  \newaliascnt{corollary}{theorem}
  \newtheorem{corollary}[corollary]{\notescorollaryname}
  \aliascntresetthe{corollary}

  \theoremstyle{definition}
  \newaliascnt{definition}{theorem}
  \newtheorem{definition}[definition]{\notesdefinitionname}
  \aliascntresetthe{definition}
  \newaliascnt{example}{theorem}
  \newtheorem{example}[example]{\notesexamplename}
  \aliascntresetthe{example}
  \newaliascnt{problem}{theorem}
  \newtheorem{problem}[problem]{\notesproblemname}
  \aliascntresetthe{problem}
  \newtheorem{exercise}{\notesexercisename}

  \theoremstyle{remark}
  \newaliascnt{remark}{theorem}
  \newtheorem{remark}[remark]{\notesremarkname}
  \aliascntresetthe{remark}

  \newenvironment{solution}
  {
  \begin{proof}[\notessolutionname]}
    {
  \end{proof}}

  \usepackage[nameinlink,noabbrev]{cleveref}

  \crefname{theorem}{theorem}{theorems}
  \crefname{lemma}{lemma}{lemmas}
  \crefname{proposition}{proposition}{propositions}
  \crefname{corollary}{corollary}{corollaries}
  \crefname{definition}{definition}{definitions}
  \crefname{example}{example}{examples}
  \crefname{problem}{problem}{problems}
  \crefname{exercise}{exercise}{exercises}
  \crefname{remark}{remark}{remarks}

  \ifx\notelanguage\notesfinnishlanguage
  \crefname{theorem}{lause}{lauseet}
  \crefname{lemma}{lemma}{lemmat}
  \crefname{proposition}{propositio}{propositiot}
  \crefname{corollary}{korollaari}{korollaarit}
  \crefname{definition}{määritelmä}{määritelmät}
  \crefname{example}{esimerkki}{esimerkit}
  \crefname{problem}{tehtävä}{tehtävät}
  \crefname{exercise}{tehtävä}{tehtävät}
  \crefname{remark}{huomautus}{huomautukset}
  \fi

  \newcommand{\NN}{\mathbb{N}}
  \newcommand{\NNone}{\mathbb{N}_1}
  \newcommand{\ZZ}{\mathbb{Z}}
  \newcommand{\QQ}{\mathbb{Q}}
  \newcommand{\RR}{\mathbb{R}}
  \newcommand{\CC}{\mathbb{C}}
  \newcommand{\PP}{\mathbb{P}}

  \DeclareMathOperator{\im}{im}
  \DeclareMathOperator{\rank}{rank}
  \DeclareMathOperator{\syt}{syt}
  \DeclareMathOperator{\pyj}{pyj}
  \DeclareMathOperator{\dom}{dom}
  \DeclarePairedDelimiter{\abs}{\lvert}{\rvert}
  \DeclarePairedDelimiter{\norm}{\lVert}{\rVert}
  \newcommand{\set}[1]{\left\lbrace#1\right\rbrace}
  \newcommand{\maxset}[1]{\max\set{#1}}
  \newcommand{\minset}[1]{\min\set{#1}}
  \newcommand{\powerset}[1]{\mathcal{P}\!\left(#1\right)}

\usepackage{tikz}
\definecolor{analysisblue}{HTML}{23566B}

\title{Analyysi II: jatkuvuus ja derivaatta}
\author{Joona Piirainen}
\date{Syksy 2026 -- periodi 2}

\begin{document}

\begin{abstract}
  Luentomuistiinpanot Helsingin yliopiston kurssilta
  Analyysi II: jatkuvuus ja derivaatta (MAT11004).
  Kurssi järjestetään aikavälillä 27.10.2026--15.12.2026.
\end{abstract}

\maketitle

\tableofcontents

\section{Kurssitiedot}

Kurssin viralliset tiedot löytyvät
\href{https://studies.helsinki.fi/kurssit/opintojakso/otm-1adb0b8a-2fbd-4953-853f-8cdc614fa9ae/MAT11004}{Helsingin
yliopiston kurssisivulta}.

\section{Muistiinpanot}

\subsection{27.10.2026: jatkuvuus pisteessä}

Jatkuvuus yhdistää funktion raja-arvon ja sen arvon tarkastelupisteessä:
kun $x$ lähestyy pistettä $x_0$, funktion arvot lähestyvät arvoa $f(x_0)$.

\subsubsection*{Raja-arvo ja funktion arvo}

\begin{example}
  Tarkastellaan funktiota $f\colon\RR\to\RR$,
  \[
    f(x)=
    \begin{cases}
      x, & x\ne 1,\\
      2, & x=1.
    \end{cases}
  \]
  Vaikka raja-arvo on olemassa, se ei ole sama kuin funktion arvo:
  \[
    \lim_{x\to 1}f(x)=1\ne 2=f(1).
  \]
  Siis $f$ ei ole jatkuva pisteessä $1$.
\end{example}

\begin{center}
  \begin{tikzpicture}[scale=1.15,>=stealth]
    \draw[->,gray] (-0.35,0)--(3.1,0) node[right] {$x$};
    \draw[->,gray] (0,-0.25)--(0,2.8) node[above] {$y$};
    \draw[analysisblue,very thick] (-0.2,-0.2)--(2.5,2.5)
    node[right,black] {$y=x$};
    \draw[dashed,gray] (1,0)--(1,2)--(0,2);
    \draw[dashed,gray] (0,1)--(1,1);
    \draw[analysisblue,thick,fill=white] (1,1) circle (2.7pt);
    \fill[analysisblue] (1,2) circle (2.7pt);
    \node[below] at (1,0) {$1$};
    \node[left] at (0,1) {$1$};
    \node[left] at (0,2) {$2$};
    \node[above] at (1,2.15) {$f(1)=2$};
  \end{tikzpicture}

  \small Avoin piste kuvaa puuttuvaa arvoa $(1,1)$;
  täytetty piste on funktion arvo kohdassa $x=1$.
\end{center}

\subsubsection*{Jatkuvuuden määritelmä}

\begin{definition}[Jatkuvuus sisäpisteessä]
  Olkoon $A\subseteq\RR$, $f\colon A\to\RR$ ja $x_0\in A$.
  Oletetaan, että jollakin $r>0$ pätee
  $(x_0-r,x_0+r)\subseteq A$.
  Funktio $f$ on \emph{jatkuva pisteessä $x_0$}, jos
  \[
    \boxed{\lim_{x\to x_0}f(x)=f(x_0).}
  \]
\end{definition}

\subsubsection*{Jatkuvuuden epsilon--delta-määritelmä}

Edellä esitetyin oletuksin $f$ on jatkuva pisteessä $x_0$ täsmälleen
silloin, kun jokaista $\varepsilon>0$ kohti on olemassa $\delta>0$
siten, että kaikilla $x\in A$ pätee
\[
  \boxed{\abs{x-x_0}<\delta
    \quad\Longrightarrow\quad
  \abs{f(x)-f(x_0)}<\varepsilon.}
\]
Luku $\varepsilon$ määrää sallitun virheen funktion arvoissa;
luku $\delta$ kertoo, kuinka lähellä pistettä $x_0$ on pysyttävä.

\begin{remark}
  Valinta $\delta=\delta(\varepsilon,x_0)$ saa riippua sekä
  luvusta $\varepsilon$ että tarkastelupisteestä $x_0$.
  Sen on kuitenkin toimittava kaikille ehdon $\abs{x-x_0}<\delta$
  täyttäville $x\in A$ samanaikaisesti.
\end{remark}

\subsubsection*{Esimerkki: jatkuvuuden osoittaminen määritelmästä}

\begin{example}
  Olkoon $f\colon\RR\to\RR$,
  \[
    f(x)=\frac12x^2+x+1.
  \]
  Osoitetaan, että $f$ on jatkuva pisteessä $-3$.
\end{example}

\begin{solution}
  \textbf{1. Hajotetaan erotus tekijöihin.}
  Koska $f(-3)=\frac52$, saadaan
  \begin{align*}
    \abs{f(x)-f(-3)}
    &=\abs{\tfrac12x^2+x+1-\tfrac52}\\
    &=\tfrac12\abs{x^2+2x-3}
    =\tfrac12\abs{x+3}\abs{x-1}.
  \end{align*}
  Tekijä $\abs{x+3}$ on juuri etäisyys tarkastelupisteestä $-3$.

  \medskip
  \noindent\textbf{2. Rajoitetaan toinen tekijä.}
  Kun $\abs{x+3}<1$, kolmioepäyhtälö antaa
  \[
    \abs{x-1}=\abs{(x+3)-4}
    \le\abs{x+3}+4<5.
  \]
  Siksi $\abs{f(x)-f(-3)}\le\frac52\abs{x+3}$.

  \medskip
  \noindent\textbf{3. Valitaan delta ja tarkistetaan ehto.}
  Olkoon $\varepsilon>0$ annettu. Valitaan
  \[
    \boxed{\delta=\min\left\{1,\frac{2\varepsilon}{5}\right\}>0.}
  \]
  Jos $\abs{x-(-3)}=\abs{x+3}<\delta$, niin
  \[
    \abs{f(x)-f(-3)}
    \le\frac52\abs{x+3}
    <\frac52\delta
    \le\varepsilon.
  \]
  Näin ollen $f$ on jatkuva pisteessä $-3$.
\end{solution}

\medskip
\noindent\textbf{Deltan valinta.}
Valinnan $\delta\le 1$ avulla saadaan paikallinen yläraja toiselle
tekijälle. Ehto $\delta\le 2\varepsilon/5$ puolestaan takaa
halutun tarkkuuden funktion arvoille.

\newpage
\subsubsection*{Esimerkki: epäjatkuvuus suoraan määritelmästä}

Jatkuvuuden määritelmän negaatio antaa tavan osoittaa epäjatkuvuus:
riittää löytää \emph{yksi} $\varepsilon_0>0$, jolle \emph{jokaista}
$\delta>0$ kohti löytyy $x\in A$ siten, että
\[
  \boxed{\abs{x-x_0}<\delta
    \quad\text{ja}\quad
  \abs{f(x)-f(x_0)}\ge\varepsilon_0.}
\]

\begin{example}
  Palataan luennon alun paloittain määriteltyyn funktioon
  $f\colon\RR\to\RR$,
  \[
    f(x)=
    \begin{cases}
      x, & x\ne 1,\\
      2, & x=1.
    \end{cases}
  \]
  Osoitetaan ilman raja-arvoja, ettei $f$ ole jatkuva pisteessä $1$.
\end{example}

\begin{solution}
  \textbf{1. Kiinnitetään sallittu virhe.}
  Valitaan $\varepsilon_0=\frac12$.

  \medskip
  \noindent\textbf{2. Etsitään vastaesimerkki jokaiselle deltalle.}
  Olkoon $\delta>0$ mielivaltainen. Valitaan
  \[
    x=1-\frac{\delta}{2}.
  \]
  Tällöin $x<1$, joten $f(x)=x$, ja piste on vaaditun lähellä:
  \[
    \abs{x-1}=\frac{\delta}{2}<\delta.
  \]
  Funktion arvot eivät kuitenkaan ole riittävän lähellä toisiaan:
  \[
    \abs{f(x)-f(1)}
    =\abs{1-\frac{\delta}{2}-2}
    =1+\frac{\delta}{2}
    >1>\varepsilon_0.
  \]
  Siis valitulla $\varepsilon_0$ mikään $\delta>0$ ei toteuta
  jatkuvuuden ehtoa. Näin ollen $f$ ei ole jatkuva pisteessä $1$.
\end{solution}

\medskip
\noindent\textbf{Kvanttorien järjestys.}
Ensin kiinnitetään $\varepsilon_0$. Sen jälkeen tarkastellaan
mielivaltaista $\delta>0$ ja valitaan siitä riippuva vastaesimerkkipiste
$x$. Yhden deltan epäonnistuminen ei vielä todista epäjatkuvuutta;
ehdon on epäonnistuttava jokaisella deltalla.

\newpage
\subsection{29.10.2026: toispuoleinen jatkuvuus ja jatkuvuuden laskusäännöt}

\subsubsection*{Kertaus ja esimerkki epäjatkuvuudesta}

Sisäpisteessä $x_0$ jatkuvuus tarkoittaa ehtoa
\[
  \lim_{x\to x_0}f(x)=f(x_0).
\]
Epsilon--delta-muodossa jokaista
$\varepsilon>0$ kohti on löydettävä $\delta>0$, jolla
\[
  x\in A,\quad \abs{x-x_0}<\delta
  \quad\Longrightarrow\quad \abs{f(x)-f(x_0)}<\varepsilon.
\]
Epäjatkuvuuden osoittamiseksi etsitään yksi $\varepsilon_0>0$, jolle
mikään $\delta>0$ ei toimi.

\begin{example}[Hyppyfunktio]
  Olkoon $f\colon\RR\to\RR$,
  \[
    f(x)=
    \begin{cases}0,&x<0,\\1,&x\ge0.
    \end{cases}
  \]
  Funktio ei ole jatkuva pisteessä $0$.
\end{example}
\begin{proof}
  Valitaan $\varepsilon_0=\frac12$. Olkoon $\delta>0$ mielivaltainen
  ja asetetaan $x=-\delta/2$. Tällöin
  \[
    \abs{x-0}=\frac{\delta}{2}<\delta,
    \qquad \abs{f(x)-f(0)}=\abs{0-1}=1>\varepsilon_0.
  \]
  Jatkuvuuden ehto ei siis toteudu millään deltalla.
\end{proof}

\subsubsection*{Perusfunktioiden jatkuvuus määritelmästä}

\begin{example}[Vakiofunktio ja lineaarinen funktio]
  Vakiofunktio $f(x)=c$ on jatkuva jokaisessa $x_0\in\RR$:
  kun $\varepsilon>0$, voidaan valita esimerkiksi $\delta=1$, sillä
  \[
    \abs{f(x)-f(x_0)}=\abs{c-c}=0<\varepsilon.
  \]
  Jos $f(x)=cx$ ja $c\ne0$, valitaan
  $\delta=\varepsilon/\abs{c}>0$. Kun $\abs{x-x_0}<\delta$, saadaan
  \[
    \abs{f(x)-f(x_0)}=\abs{c}\abs{x-x_0}
    <\abs{c}\delta=\varepsilon.
  \]
  Tapaus $c=0$ on jo käsitelty vakiofunktiona. Erityisesti
  identtinen funktio $x\mapsto x$ on jatkuva.
\end{example}

\begin{example}[Kuutiofunktio avoimella välillä]
  Funktio $f\colon(0,1)\to\RR$, $f(x)=x^3$, on jatkuva.
  Kun $x,x_0\in(0,1)$, erotuksen tekijöihinjako antaa
  \[
    \abs{x^3-x_0^3}
    =\abs{x-x_0}\abs{x^2+xx_0+x_0^2}
    \le3\abs{x-x_0}.
  \]
  Siksi jokaiselle $\varepsilon>0$ riittää
  $\delta=\varepsilon/3$: jos $x\in(0,1)$ ja
  $\abs{x-x_0}<\delta$, niin $\abs{x^3-x_0^3}<\varepsilon$.
  Määrittelyjoukkoon kuulumista ei tarvitse varmistaa deltalla,
  sillä $x$ valitaan alun perin tästä joukosta.
\end{example}

\newpage
\subsubsection*{Vasemmalta ja oikealta jatkuvuus}

\begin{definition}[Toispuoleinen jatkuvuus]
  Olkoot $f\colon A\to\RR$ ja $x_0\in A$.
  Jos $(x_0-r,x_0]\subseteq A$ jollakin $r>0$, niin $f$ on
  \emph{vasemmalta jatkuva} pisteessä $x_0$, kun
  \[
    \lim_{x\to x_0^-}f(x)=f(x_0).
  \]
  Jos $[x_0,x_0+r)\subseteq A$ jollakin $r>0$, niin $f$ on
  \emph{oikealta jatkuva} pisteessä $x_0$, kun
  \[
    \lim_{x\to x_0^+}f(x)=f(x_0).
  \]
\end{definition}

Oikealta jatkuvuuden epsilon--delta-ehto on: jokaista
$\varepsilon>0$ kohti on $\delta>0$ siten, että kaikilla $x\in A$
\[
  0<x-x_0<\delta\quad\Longrightarrow\quad
  \abs{f(x)-f(x_0)}<\varepsilon.
\]
Vasemmalta jatkuvuudessa lähtöehto on $0<x_0-x<\delta$.
Pisteen $x=x_0$ voi myös sisällyttää ehtoihin, koska silloin
funktion arvojen erotus on nolla.

\begin{example}
  Olkoon
  \[
    f(x)=
    \begin{cases}x^2,&x<0,\\x^2+1,&x\ge0.
    \end{cases}
  \]
  Tällöin $f(0)=1$. Funktio on oikealta jatkuva mutta ei
  vasemmalta jatkuva pisteessä $0$.
\end{example}
\begin{solution}
  Kun $\varepsilon>0$, valinta $\delta=\sqrt{\varepsilon}$ antaa
  $0<x<\delta\Rightarrow\abs{f(x)-f(0)}=x^2<\varepsilon$.
  Siis $f$ on oikealta jatkuva. Samalla arviolla
  $-\delta<x<0\Rightarrow\abs{f(x)-0}=x^2<\varepsilon$, joten
  vasemmanpuoleinen raja-arvo on $0\ne1=f(0)$.
  Funktio ei siis ole vasemmalta jatkuva.
\end{solution}

\begin{theorem}[Jatkuvuus ja toispuoleinen jatkuvuus]
  Oletetaan, että $(x_0-r,x_0+r)\subseteq A$ jollakin $r>0$.
  Funktio $f\colon A\to\RR$ on jatkuva pisteessä $x_0$ täsmälleen
  silloin, kun se on sekä vasemmalta että oikealta jatkuva siinä.
\end{theorem}
\begin{proof}
  Jos $f$ on jatkuva pisteessä $x_0$, jatkuvuuden ehdon toteuttava
  $\delta$ toimii myös silloin, kun $x$ rajoitetaan kummalle tahansa
  puolelle pistettä $x_0$. Siis $f$ on jatkuva molemmilta puolilta.

  Oletetaan kääntäen toispuoleiset jatkuvuudet ja kiinnitetään
  $\varepsilon>0$. Valitaan niiden perusteella $\delta_->0$ ja
  $\delta_+>0$ siten, että $\abs{f(x)-f(x_0)}<\varepsilon$, kun
  $x_0-\delta_-<x<x_0$ tai $x_0<x<x_0+\delta_+$.
  Asetetaan $\delta=\min\{\delta_-,\delta_+\}$.
  Jos $\abs{x-x_0}<\delta$ ja $x\ne x_0$, toinen näistä ehdoista
  toteutuu. Jos $x=x_0$, erotus on nolla. Näin $\delta$ toteuttaa
  kaksipuolisen jatkuvuuden ehdon.
\end{proof}

\newpage
\subsubsection*{Jatkuvuus joukossa ja välin päätepisteissä}

\begin{definition}[Jatkuvuus välillä]
  Funktio on \emph{jatkuva avoimella välillä} $(a,b)$, jos se on
  jatkuva jokaisessa välin pisteessä. Funktio on
  \emph{jatkuva suljetulla välillä} $[a,b]$, jos se on
  \begin{itemize}
    \item jatkuva jokaisessa sisäpisteessä $x\in(a,b)$,
    \item oikealta jatkuva pisteessä $a$ ja
    \item vasemmalta jatkuva pisteessä $b$.
  \end{itemize}
  Puoliavoimella välillä tarkistetaan vastaavasti vain mukana
  olevan päätepisteen jatkuvuus välin sisältä päin.
\end{definition}

\begin{definition}[Jatkuvuus joukossa]
  Funktio $f\colon A\to\RR$ on \emph{jatkuva}, jos jokaista
  $x_0\in A$ ja $\varepsilon>0$ kohti on olemassa $\delta>0$, jolla
  \[
    x\in A,\quad\abs{x-x_0}<\delta
    \quad\Longrightarrow\quad\abs{f(x)-f(x_0)}<\varepsilon.
  \]
  Jatkuvuus osajoukossa $E\subseteq A$ tarkoittaa rajoittuman
  $f|_E\colon E\to\RR$ jatkuvuutta: tällöin sekä $x_0$ että $x$
  valitaan joukosta $E$.
\end{definition}

Tämä määritelmä toimii myös muille joukoille kuin väleille.
Välillä se antaa juuri edellä kuvatut ehdot: sisäpistettä voidaan
lähestyä molemmilta puolilta, mutta päätepistettä vain välin sisältä.
Suljetun välin vasemmassa päätepisteessä $a$ ehto
$x\in[a,b]$ ja $\abs{x-a}<\delta$ tarkoittaa $a\le x<a+\delta$;
oikeassa päätepisteessä saadaan vastaavasti $b-\delta<x\le b$.

\begin{example}[Kuutiofunktio puoliavoimella välillä]
  Funktio $f\colon[0,1)\to\RR$, $f(x)=x^3$, on jatkuva.
  Edellä saatu arvio $\abs{x^3-x_0^3}\le3\abs{x-x_0}$ pätee
  myös kaikille $x,x_0\in[0,1)$. Valinta $\delta=\varepsilon/3$
  todistaa siis jatkuvuuden samalla kertaa sekä sisäpisteissä
  että päätepisteessä $0$. Piste $1$ ei kuulu määrittelyjoukkoon.
\end{example}

\begin{example}[Määrittelyjoukon merkitys]
  Olkoon $f\colon\RR\setminus\{0\}\to\RR$,
  \[
    f(x)=
    \begin{cases}-1,&x<0,\\1,&x>0.
    \end{cases}
  \]
  Funktio on jatkuva. Nimittäin kun $x_0\ne0$ ja
  $\varepsilon>0$, valitaan $\delta=\abs{x_0}>0$.
  Jos $\abs{x-x_0}<\delta$, niin $x$ ja $x_0$ ovat samanmerkkisiä,
  joten
  \[
    \abs{f(x)-f(x_0)}=0<\varepsilon.
  \]
  Kuvaajan kaksi erillistä osaa eivät estä jatkuvuutta:
  $0$ ei kuulu funktion määrittelyjoukkoon.
\end{example}

\newpage
\subsubsection*{Summien, tulojen ja osamäärien jatkuvuus}

\begin{theorem}[Jatkuvuuden laskusäännöt]
  Olkoot $f,g\colon A\to\RR$ jatkuvia sisäpisteessä $x_0$.
  Tällöin $f+g$ ja $fg$ ovat jatkuvia pisteessä $x_0$.
  Jos lisäksi $g(x_0)\ne0$, niin myös $f/g$ on jatkuva
  pisteessä $x_0$.
\end{theorem}
\begin{proof}
  Merkitään $a=f(x_0)$ ja $b=g(x_0)$ ja kiinnitetään
  $\varepsilon>0$. Kaikissa seuraavissa arvioissa $x\in A$.

  \smallskip
  \noindent\textbf{Summa.}
  Jatkuvuuksien nojalla on $\delta_f,\delta_g>0$ siten, että
  \[
    \abs{x-x_0}<\delta_f\Rightarrow\abs{f(x)-a}<\varepsilon/2,
    \qquad
    \abs{x-x_0}<\delta_g\Rightarrow\abs{g(x)-b}<\varepsilon/2.
  \]
  Kun $\abs{x-x_0}<\min\{\delta_f,\delta_g\}$, kolmioepäyhtälö antaa
  \[
    \abs{(f(x)+g(x))-(a+b)}
    \le\abs{f(x)-a}+\abs{g(x)-b}<\varepsilon.
  \]

  \smallskip
  \noindent\textbf{Tulo.}
  Jatkuvuuksien nojalla voidaan valita yhteinen $\delta>0$ siten,
  että ehdosta $\abs{x-x_0}<\delta$ seuraa
  \[
    \abs{f(x)-a}<\min\left\{1,\frac{\varepsilon}{2(\abs{b}+1)}\right\},
    \qquad
    \abs{g(x)-b}<\frac{\varepsilon}{2(\abs{a}+1)}.
  \]
  Tällöin $\abs{f(x)}\le\abs{a}+\abs{f(x)-a}<\abs{a}+1$ ja
  \begin{align*}
    \abs{f(x)g(x)-ab}
    &=\abs{f(x)(g(x)-b)+b(f(x)-a)}\\
    &\le\abs{f(x)}\abs{g(x)-b}+\abs{b}\abs{f(x)-a}\\
    &<\frac{\varepsilon}{2}+\frac{\varepsilon}{2}=\varepsilon.
  \end{align*}

  \smallskip
  \noindent\textbf{Osamäärä.}
  Oletetaan $b\ne0$. Osoitetaan ensin funktion $1/g$ jatkuvuus.
  Valitaan $\delta>0$ siten, että ehdosta $\abs{x-x_0}<\delta$ seuraa
  \[
    \abs{g(x)-b}<\min\left\{\frac{\abs{b}}2,
    \frac{\varepsilon\abs{b}^2}{2}\right\}.
  \]
  Käänteinen kolmioepäyhtälö antaa
  $\abs{g(x)}\ge\abs{b}-\abs{g(x)-b}>\abs{b}/2>0$, joten
  $1/g$ on määritelty tässä ympäristössä ja
  \[
    \left|\frac1{g(x)}-\frac1b\right|
    =\frac{\abs{g(x)-b}}{\abs{g(x)}\abs{b}}
    \le\frac{2\abs{g(x)-b}}{\abs{b}^2}<\varepsilon.
  \]
  Siis $1/g$ on jatkuva pisteessä $x_0$, ja jo todistetun
  tulosäännön mukaan myös $f/g=f\cdot(1/g)$ on jatkuva siinä.
\end{proof}

Erityisesti $-f=(-1)f$ on jatkuva, jos $f$ on jatkuva.
Samoin jatkuvien funktioiden erotukset, vakiokerrannaiset sekä
äärelliset summat ja tulot ovat jatkuvia.

\newpage
\subsubsection*{Potenssit ja polynomit}

\begin{theorem}[Potenssifunktioiden jatkuvuus]
  Funktio $x\mapsto x^n$ on jatkuva koko reaalilukujoukossa
  jokaisella positiivisella kokonaisluvulla $n$.
\end{theorem}
\begin{proof}
  Todistetaan induktiolla. Kun $n=1$, väite on edellä osoitettu
  identtisen funktion jatkuvuus. Jos $x\mapsto x^n$ on jatkuva,
  niin $x\mapsto x^{n+1}=x\cdot x^n$ on kahden jatkuvan funktion
  tulona jatkuva. Induktioperiaate antaa väitteen kaikille $n\ge1$.
  Myös vakiopolynomi $x\mapsto1$ on jatkuva.
\end{proof}

\begin{theorem}[Polynomien jatkuvuus]
  Jokainen polynomifunktio $P\colon\RR\to\RR$ on jatkuva.
\end{theorem}
\begin{proof}
  Kirjoitetaan
  \[
    P(x)=a_nx^n+a_{n-1}x^{n-1}+\cdots+a_1x+a_0.
  \]
  Potenssifunktiot ja vakiofunktiot ovat jatkuvia.
  Siis jokainen termi on jatkuva, ja äärellisen summan jatkuvuus
  antaa väitteen.
\end{proof}

\begin{corollary}
  Rationaalifunktio $P/Q$, missä $P$ ja $Q$ ovat polynomeja,
  on jatkuva joukossa $\{x\in\RR:Q(x)\ne0\}$.
\end{corollary}
\begin{proof}
  Olkoon $x_0$ mikä tahansa nimittäjän nollakohdista eroava piste.
  Polynomit $P$ ja $Q$ ovat jatkuvia pisteessä $x_0$, ja
  $Q(x_0)\ne0$. Osamääräsäännön mukaan $P/Q$ on jatkuva siinä.
  Koska $x_0$ oli mielivaltainen, väite pätee koko määrittelyjoukossa.
\end{proof}

\medskip
\noindent\textbf{Todistusten työnjako.}
Epsilon--delta-määritelmällä rakennetaan perustulokset.
Laskusäännöillä niistä saadaan laajoja funktioluokkia ilman,
että jokaiselle lausekkeelle etsitään erikseen delta.

\subsubsection*{Yhdistetyn funktion jatkuvuus}

\begin{theorem}[Yhdistetyn funktion jatkuvuus]
  Olkoot $A,B\subseteq\RR$, $f\colon A\to B$ ja $g\colon B\to\RR$.
  Jos $f$ on jatkuva pisteessä $x_0\in A$ ja $g$ on jatkuva
  pisteessä $f(x_0)$, niin $g\circ f$ on jatkuva pisteessä $x_0$.
\end{theorem}
\begin{proof}
  Olkoon $\varepsilon>0$.
  Koska $g$ on jatkuva pisteessä $f(x_0)$, on olemassa $\eta>0$
  siten, että kaikilla $y\in B$
  \[
    \abs{y-f(x_0)}<\eta
    \quad\Longrightarrow\quad
    \abs{g(y)-g(f(x_0))}<\varepsilon.
  \]
  Koska $f$ on jatkuva pisteessä $x_0$, tätä lukua $\eta$ kohti
  on olemassa $\delta>0$ siten, että kaikilla $x\in A$
  \[
    \abs{x-x_0}<\delta
    \quad\Longrightarrow\quad
    \abs{f(x)-f(x_0)}<\eta.
  \]
  Asettamalla ensimmäiseen ehtoon $y=f(x)$ saadaan
  \[
    \abs{x-x_0}<\delta
    \quad\Longrightarrow\quad
    \abs{g(f(x))-g(f(x_0))}<\varepsilon.
  \]
  Tämä on funktion $g\circ f$ jatkuvuuden ehto.
\end{proof}

Ulomman funktion jatkuvuus määrää ensin välivaiheen tarkkuuden
$\eta$; sisemmän funktion jatkuvuus antaa sitten sitä vastaavan $\delta$:n.

\begin{example}[Laskusäännöt ja yhdistäminen yhdessä]
  Funktio
  \[
    F\colon\RR\to\RR,\qquad
    F(x)=\left(\frac{5x^4+3x+1}{(x+1)^2+1}\right)^4
  \]
  on jatkuva.
\end{example}
\begin{solution}
  Osoittaja ja nimittäjä ovat polynomeja, ja
  $(x+1)^2+1\ge1>0$ kaikilla $x$. Niiden osamäärä on siis jatkuva
  koko joukossa $\RR$. Yhdistäminen jatkuvan funktion $t\mapsto t^4$
  kanssa säilyttää jatkuvuuden.
\end{solution}

\newpage
\subsubsection*{Paloittain jatkuvuus}

\begin{definition}[Paloittain jatkuva funktio]
  Funktio $f\colon(a,b)\to\RR$ on \emph{paloittain jatkuva},
  jos on olemassa äärellinen joukko
  $\{x_1,\ldots,x_n\}\subset(a,b)$ siten, että
  \begin{itemize}
    \item $f$ on jatkuva joukossa $(a,b)\setminus\{x_1,\ldots,x_n\}$ ja
    \item molemmat toispuoleiset raja-arvot
      $\lim_{x\to x_i^-}f(x)$ ja $\lim_{x\to x_i^+}f(x)$
      ovat olemassa ja äärellisiä jokaisella $i$.
  \end{itemize}
  Suljetulla välillä vaaditaan vastaavasti päätepisteissä
  äärelliset raja-arvot välin sisältä päin.
\end{definition}

\begin{remark}
  Raja-arvojen ei tarvitse olla keskenään samat eikä niiden tarvitse
  olla yhtä suuria kuin $f(x_i)$. Äärellinen määrä hyppyjä tai
  poistuvia epäjatkuvuuksia on siis sallittu. Pelkkä jatkuvuus
  äärellisen joukon ulkopuolella ei riitä: myös toispuoleisten
  raja-arvojen olemassaolo ja äärellisyys on tarkistettava.
\end{remark}

\begin{example}
  Olkoon $f\colon(-2,2)\to\RR$,
  \[
    f(x)=
    \begin{cases}
      -1,&-2<x<-1,\\
      1,&-1\le x\le1,\\
      x+1,&1<x<2.
    \end{cases}
  \]
  Funktio on paloittain jatkuva.
\end{example}
\begin{solution}
  Väleillä $(-2,-1)$ ja $(-1,1)$ funktio on vakio, ja välillä
  $(1,2)$ se on polynomi $x+1$. Siis se on jatkuva joukon
  $\{-1,1\}$ ulkopuolella. Liitoskohtien tiedot ovat
  \[
    \begin{array}{c|ccc}
      x_i & \displaystyle\lim_{x\to x_i^-}f(x)
      & f(x_i) & \displaystyle\lim_{x\to x_i^+}f(x)\\[2pt]
      \hline
      -1 & -1 & 1 & 1\\
      1 &  1 & 1 & 2
    \end{array}
  \]
  Kaikki neljä toispuoleista raja-arvoa ovat äärellisiä, joten
  $f$ on paloittain jatkuva. Kummassakaan liitoskohdassa $f$ ei
  ole jatkuva, koska vasen ja oikea raja-arvo eroavat toisistaan.
  Se on kuitenkin oikealta jatkuva pisteessä $-1$ ja vasemmalta
  jatkuva pisteessä $1$.
\end{solution}

\newpage
\subsection{Bolzanon lause ja väliarvot}

Jatkuva funktio ei voi vaihtaa merkkiä kulkematta nollan kautta.
Tämän täsmällinen muoto on Bolzanon lause. Todistuksissa tarvitaan
reaalilukujen täydellisyyttä sekä seuraavia jatkuvuuden seurauksia.

\subsubsection*{Todistuksissa tarvittavat aputulokset}

\begin{lemma}[Merkin säilyminen lähiympäristössä]
  \label{lem:local-sign}
  Olkoon $f\colon A\to\RR$ jatkuva pisteessä $x_0\in A$.
  Jos $f(x_0)>0$, niin jollakin $\delta>0$ pätee
  \[
    x\in A,\quad\abs{x-x_0}<\delta
    \quad\Longrightarrow\quad f(x)>\tfrac12f(x_0)>0.
  \]
  Jos $f(x_0)<0$, saadaan vastaavasti
  $f(x)<\tfrac12f(x_0)<0$ riittävän lähellä pistettä $x_0$.
\end{lemma}
\begin{proof}
  Valitaan jatkuvuuden määritelmässä
  $\varepsilon=\abs{f(x_0)}/2>0$. Tällöin sopivalla $\delta>0$
  \[
    f(x_0)-\frac{\abs{f(x_0)}}2
    <f(x)<f(x_0)+\frac{\abs{f(x_0)}}2,
  \]
  kun $x\in A$ ja $\abs{x-x_0}<\delta$.
  Molemmat väitteet seuraavat tästä.
\end{proof}

\begin{theorem}[Jatkuvuuden jonokriteeri]
  \label{thm:sequential-continuity}
  Funktio $f\colon A\to\RR$ on jatkuva pisteessä $x_0\in A$
  täsmälleen silloin, kun jokaiselle jonolle $(x_n)$ joukossa $A$
  pätee
  \[
    x_n\to x_0\quad\Longrightarrow\quad f(x_n)\to f(x_0).
  \]
\end{theorem}
\begin{proof}
  Oletetaan ensin jatkuvuus ja olkoon $x_n\to x_0$.
  Kun $\varepsilon>0$, jatkuvuus antaa $\delta>0$, jolla
  $\abs{x-x_0}<\delta$ ja $x\in A$ takaavat
  $\abs{f(x)-f(x_0)}<\varepsilon$.
  Jonon suppenemisen nojalla $\abs{x_n-x_0}<\delta$ kaikilla
  riittävän suurilla $n$, joten myös
  $\abs{f(x_n)-f(x_0)}<\varepsilon$ lopulta.

  Jos $f$ ei ole jatkuva pisteessä $x_0$, on olemassa
  $\varepsilon_0>0$ siten, että jokaista $n\ge1$ kohti voidaan
  valita $x_n\in A$, jolle
  \[
    \abs{x_n-x_0}<\frac1n,
    \qquad \abs{f(x_n)-f(x_0)}\ge\varepsilon_0.
  \]
  Tällöin $x_n\to x_0$, mutta $f(x_n)$ ei suppene kohti
  lukua $f(x_0)$. Tämä todistaa käänteisen suunnan.
\end{proof}

\begin{lemma}[Sisäkkäiset suljetut välit]
  \label{lem:nested-intervals}
  Olkoot $I_n=[a_n,b_n]$ epätyhjiä suljettuja välejä, joille
  $I_{n+1}\subseteq I_n$ kaikilla $n\ge0$ ja $b_n-a_n\to0$.
  Tällöin on olemassa täsmälleen yksi
  \[
    c\in\bigcap_{n=0}^{\infty}I_n.
  \]
  Lisäksi $a_n\to c$ ja $b_n\to c$.
\end{lemma}
\begin{proof}
  Jono $(a_n)$ on kasvava ja ylhäältä rajoitettu luvulla $b_0$.
  Olkoon $c=\sup\{a_n:n\ge0\}$. Jokainen $b_n$ on kaikkien
  lukujen $a_k$ yläraja: kun $k\le n$, niin $a_k\le a_n\le b_n$,
  ja kun $k>n$, sisäkkäisyys antaa $a_k\le b_k\le b_n$.
  Siis $a_n\le c\le b_n$ kaikilla $n$.
  Koska
  \[
    0\le c-a_n\le b_n-a_n,
    \qquad 0\le b_n-c\le b_n-a_n,
  \]
  päätepistejonot suppenevat kohti lukua $c$.
  Jos myös $d$ kuuluu kaikkiin väleihin, niin
  $\abs{c-d}\le b_n-a_n$ kaikilla $n$, joten $c=d$.
\end{proof}

\newpage
\subsubsection*{Bolzanon lause: todistus supremumin avulla}

\begin{theorem}[Bolzanon lause]
  \label{thm:bolzano}
  Olkoot $a<b$ ja $f\colon[a,b]\to\RR$ jatkuva.
  Jos $f(a)f(b)<0$, niin on olemassa $c\in(a,b)$ siten, että
  \[
    \boxed{f(c)=0.}
  \]
\end{theorem}
\begin{proof}
  Oletetaan ensin $f(a)<0<f(b)$ ja asetetaan
  \[
    E=\{x\in[a,b]:f(x)<0\}.
  \]
  Joukko on epätyhjä, sillä $a\in E$, ja ylhäältä rajoitettu,
  sillä $E\subseteq[a,b]$. Reaalilukujen täydellisyyden nojalla
  on olemassa $s=\sup E\in[a,b]$.

  Osoitetaan ensin $a<s<b$.
  Oikealta jatkuvuus pisteessä $a$ ja \cref{lem:local-sign}
  antavat $\delta_a>0$, jolla $f(x)<0$, kun
  $x\in[a,b]$ ja $x<a+\delta_a$. Esimerkiksi
  \[
    t=a+\min\{\delta_a/2,(b-a)/2\}
  \]
  kuuluu joukkoon $E$ ja on suurempi kuin $a$.
  Siis $s\ge t>a$.
  Vastaavasti vasemmalta jatkuvuus pisteessä $b$ antaa
  $\delta_b>0$, jolla $f(x)>0$, kun
  $x\in[a,b]$ ja $b-\delta_b<x\le b$.
  Kun $h=\min\{\delta_b/2,(b-a)/2\}>0$, kaikki joukon $E$
  alkiot ovat enintään $b-h$. Siten $s\le b-h<b$.

  Supremumin määritelmän nojalla voidaan valita jokaisella
  $n\ge1$ piste $x_n\in E$, jolle
  \[
    s-\frac1n<x_n\le s.
  \]
  Tällöin $x_n\to s$. Jatkuvuuden jonokriteerin mukaan
  \[
    f(s)=\lim_{n\to\infty}f(x_n)\le0,
  \]
  koska $f(x_n)<0$ kaikilla $n$. Huomaa, että raja-arvo
  voi olla nolla, vaikka jokainen jonon termi on negatiivinen.

  Toisaalta jokainen $x\in(s,b]$ on joukon $E$ ulkopuolella,
  joten $f(x)\ge0$. Valitaan esimerkiksi
  $y_n=s+(b-s)/(n+1)$. Tällöin $y_n\in(s,b]$, $y_n\to s$ ja
  \[
    f(s)=\lim_{n\to\infty}f(y_n)\ge0.
  \]
  Näin $f(s)=0$, joten voidaan valita $c=s\in(a,b)$.

  Jos $f(a)>0>f(b)$, sovelletaan jo todistettua tapausta
  jatkuvaan funktioon $-f$. Sen nollakohdat ovat samat kuin
  funktion $f$ nollakohdat.
\end{proof}

\begin{remark}
  Merkintä $a\in E$ koskee määrittelyjoukon pistettä $a$;
  funktion arvo $f(a)$ ei välttämättä kuulu joukkoon $E$.
  Jos päätepisteessä on jo $f(a)=0$ tai $f(b)=0$, nollakohta
  löytyy päätepisteestä. Ehto $f(a)f(b)<0$ takaa nollakohdan
  nimenomaan avoimella välillä $(a,b)$.
\end{remark}

\newpage
\subsubsection*{Toinen todistus: välin puolittaminen}

\begin{proof}[Bolzanon lauseen toinen todistus]
  Oletetaan $f(a)>0>f(b)$ ja asetetaan $I_0=[a,b]$.
  Muodostetaan induktiivisesti välit $I_n=[a_n,b_n]$.
  Kun $I_n$ on valittu, olkoon
  \[
    c_n=\frac{a_n+b_n}{2},\qquad
    I_{n+1}=
    \begin{cases}
      [c_n,b_n],&f(c_n)\ge0,\\
      [a_n,c_n],&f(c_n)<0.
    \end{cases}
  \]
  Jos $f(c_n)=0$, nollakohta on jo löytynyt; määritelmä sallii
  myös puolittamisen jatkamisen. Jokaisella $n\ge0$ pätee
  \[
    I_{n+1}\subseteq I_n,\qquad
    f(a_n)\ge0,\qquad f(b_n)<0,\qquad
    b_n-a_n=\frac{b-a}{2^n}\longrightarrow0.
  \]
  \Cref{lem:nested-intervals} antaa täsmälleen yhden pisteen
  $c\in\bigcap_{n=0}^{\infty}I_n$ sekä suppenemiset
  $a_n\to c$ ja $b_n\to c$. Jatkuvuuden nojalla
  \[
    f(c)=\lim_{n\to\infty}f(a_n)\ge0,
    \qquad
    f(c)=\lim_{n\to\infty}f(b_n)\le0.
  \]
  Siis $f(c)=0$. Koska alkuperäiset päätepistearvot ovat
  nollasta poikkeavia, $c\in(a,b)$.
  Toinen merkkijärjestys käsitellään korvaamalla $f$ funktiolla $-f$.
\end{proof}

\begin{remark}[Puolitusmenetelmän virhearvio]
  Kun nollakohta $c$ kuuluu väliin $I_n$, sen keskipiste $c_n$
  toteuttaa
  \[
    \abs{c_n-c}\le\frac{b_n-a_n}{2}
    =\frac{b-a}{2^{n+1}}.
  \]
  Todistus antaa siis myös numeerisen menetelmän nollakohdan
  arvioimiseen. Lause takaa ainakin yhden nollakohdan;
  nollakohdan yksikäsitteisyys vaatii lisäoletuksen,
  esimerkiksi aidon monotonisuuden.
\end{remark}

\subsubsection*{Väliarvolause}

\begin{corollary}[Väliarvolause]
  \label{cor:intermediate-value}
  Jos $f\colon[a,b]\to\RR$ on jatkuva, se saa jokaisen arvon
  lukujen $f(a)$ ja $f(b)$ välissä. Täsmällisesti
  \[
    \min\{f(a),f(b)\}\le y\le\max\{f(a),f(b)\}
    \quad\Longrightarrow\quad
    \text{on olemassa }c\in[a,b]\text{, jolla }f(c)=y.
  \]
  Jos epäyhtälöt ovat aitoja, voidaan valita $c\in(a,b)$.
\end{corollary}
\begin{proof}
  Päätepistearvot saadaan valitsemalla $c=a$ tai $c=b$.
  Jos $y$ on aidosti niiden välissä, funktio
  $g(x)=f(x)-y$ on jatkuva ja sen päätepistearvot ovat
  erimerkkiset. Bolzanon lause antaa $c\in(a,b)$, jolla
  $g(c)=0$, eli $f(c)=y$.
\end{proof}

\newpage
\subsubsection*{Jatkuva kuva välistä on väli}

\begin{lemma}
  \label{lem:interval-image}
  Olkoon $\Delta\subseteq\RR$ epätyhjä väli ja
  $f\colon\Delta\to\RR$ jatkuva. Tällöin
  \[
    f(\Delta)=\{f(x):x\in\Delta\}
  \]
  on väli; myös yhden pisteen joukko luetaan väliksi.
\end{lemma}
\begin{proof}
  Olkoot $y_1,y_2\in f(\Delta)$ ja $y_1<y_2$.
  Valitaan $x_1,x_2\in\Delta$ siten, että
  $f(x_1)=y_1$ ja $f(x_2)=y_2$.
  Koska $\Delta$ on väli,
  \[
    A=[\min\{x_1,x_2\},\max\{x_1,x_2\}]\subseteq\Delta.
  \]
  Funktion $f$ rajoittuma välille $A$ on jatkuva.
  Väliarvolauseen mukaan jokaista $y\in(y_1,y_2)$ kohti on
  $x\in A$, jolla $f(x)=y$. Siis $y\in f(\Delta)$.
  Kuva sisältää näin kaikkien kahden alkionsa väliset reaaliluvut,
  mikä on välin tunnusomainen ominaisuus.
\end{proof}

\begin{proof}[Vaihtoehtoinen todistus infimumin ja supremumin avulla]
  Merkitään laajennettuja reaalilukuja käyttäen
  \[
    p=\inf f(\Delta)\in\RR\cup\{-\infty\},\qquad
    s=\sup f(\Delta)\in\RR\cup\{\infty\}.
  \]
  Jos $p=s$, funktio on vakio ja kuva on yhden pisteen joukko.
  Oletetaan $p<s$ ja valitaan reaaliluku $y$, jolla $p<y<s$.
  Infimumin ja supremumin määritelmistä seuraa, että on olemassa
  $x_1,x_2\in\Delta$ siten, että
  \[
    f(x_1)<y<f(x_2).
  \]
  Tämä pätee myös, kun $p=-\infty$ tai $s=\infty$:
  tällöin käytetään vastaavasti kuvan rajoittamattomuutta.
  Kuten edellisessä todistuksessa, väliarvolause antaa
  $x\in\Delta$, jolla $f(x)=y$.
  Näin kuva sisältää kaikki reaaliluvut aidosti rajojen $p$ ja $s$
  välissä, eikä se sisällä näiden rajojen ulkopuolisia lukuja.
  Äärellinen päätepiste kuuluu kuvaan täsmälleen silloin,
  kun funktio saavuttaa kyseisen arvon. Siis kuva on väli.
\end{proof}

\begin{remark}
  Välin kuva voi olla avoin, suljettu, puoliavoin tai rajoittamaton.
  Esimerkiksi $f(x)=x^2$ kuvaa välin $(-1,1)$ väliksi $[0,1)$.
  Lähtövälin oletus on olennainen: jatkuvan funktion
  $f(x)=x$ kuva joukosta $\RR\setminus\{0\}$ ei ole väli.
\end{remark}

\newpage
\subsection{Käänteisfunktion jatkuvuus ja juurifunktiot}

\subsubsection*{Aidosti monotonisen funktion käänteisfunktio}

\begin{theorem}[Käänteisfunktion jatkuvuus]
  \label{thm:inverse-continuity}
  Olkoon $\Delta\subseteq\RR$ epätyhjä väli ja
  $f\colon\Delta\to\RR$ jatkuva ja aidosti monotoninen.
  Tällöin käänteisfunktio
  \[
    f^{-1}\colon f(\Delta)\to\Delta
  \]
  on olemassa, jatkuva ja aidosti monotoninen samaan suuntaan
  kuin $f$. Lisäksi $f(\Delta)$ on väli.
\end{theorem}
\begin{proof}
  Aito monotonisuus takaa injektiivisyyden, ja funktio on
  surjektiivinen omaan kuvaansa. Käänteisfunktio $g=f^{-1}$
  on siis olemassa. Jos $f$ on aidosti kasvava ja $y_1<y_2$,
  niin $g(y_1)<g(y_2)$, sillä vastakkainen järjestys antaisi
  $y_1=f(g(y_1))\ge f(g(y_2))=y_2$.
  Siis myös $g$ on aidosti kasvava.

  Todistetaan jatkuvuus jonokriteerillä. Olkoon
  $y_0\in f(\Delta)$, $x_0=g(y_0)$ ja $y_n\to y_0$, missä
  $y_n\in f(\Delta)$. Oletetaan vastoin väitettä, ettei
  $g(y_n)\to x_0$. Tällöin jollakin $\varepsilon>0$ on
  osajono $(y_{n_k})$, jolle
  \[
    \abs{g(y_{n_k})-x_0}\ge\varepsilon
    \quad\text{kaikilla }k.
  \]
  Ainakin toinen erotuksen merkeistä esiintyy äärettömän usein.
  Siirtymällä edelleen osajonoon voidaan siis olettaa joko
  \[
    g(y_{n_k})\ge x_0+\varepsilon\quad\text{kaikilla }k,
    \qquad\text{tai}\qquad
    g(y_{n_k})\le x_0-\varepsilon\quad\text{kaikilla }k.
  \]
  Ensimmäisessä tapauksessa $z=x_0+\varepsilon/2\in\Delta$,
  sillä $z$ on pisteiden $x_0$ ja $g(y_{n_1})$ välissä ja
  $\Delta$ on väli. Aidon kasvavuuden nojalla
  \[
    y_{n_k}=f(g(y_{n_k}))\ge f(z)>f(x_0)=y_0.
  \]
  Tämä on ristiriidassa suppenemisen $y_{n_k}\to y_0$ kanssa,
  sillä kiinteä luku $f(z)-y_0$ on positiivinen.
  Toisessa tapauksessa $z=x_0-\varepsilon/2\in\Delta$ ja
  \[
    y_{n_k}=f(g(y_{n_k}))\le f(z)<f(x_0)=y_0,
  \]
  mikä antaa samanlaisen ristiriidan.
  Siis $g(y_n)\to g(y_0)$, joten $g$ on jatkuva.

  Jos $f$ on aidosti vähenevä, $h=-f$ on jatkuva ja aidosti
  kasvava. Jo todistetun nojalla $h^{-1}$ on jatkuva ja aidosti
  kasvava. Koska $f^{-1}(y)=h^{-1}(-y)$, myös $f^{-1}$ on
  jatkuva ja se on aidosti vähenevä.
  Väite kuvan välimuodosta on \cref{lem:interval-image}.
\end{proof}

\begin{remark}[Päätepisteet käänteisfunktion todistuksessa]
  Pisteiden $x_0\pm\varepsilon$ kuulumista välille $\Delta$
  ei voi olettaa, varsinkaan päätepisteessä.
  Edellä valittu väliin jäävä piste $z$ kuuluu varmasti välille
  siinä tapauksessa, jossa sitä käytetään.
  Jatkuvuus tarkoittaa kaikkialla jatkuvuutta suhteessa
  määrittelyjoukkoon, joten todistus kattaa myös päätepisteet.
\end{remark}

\newpage
\subsubsection*{Juurifunktion olemassaolo ja jatkuvuus}

\begin{theorem}[Positiivinen kokonaislukujuuri]
  \label{thm:nth-root}
  Jokaisella kokonaisluvulla $n\ge1$ ja luvulla $x\ge0$ on
  täsmälleen yksi $u\ge0$, jolle $u^n=x$. Tätä lukua merkitään
  \[
    u=\sqrt[n]{x}=x^{1/n}.
  \]
  Funktio $x\mapsto x^{1/n}$ on jatkuva ja aidosti kasvava
  välillä $[0,\infty)$.
\end{theorem}
\begin{proof}
  Funktio $P_n(u)=u^n$ on jatkuva ja aidosti kasvava välillä
  $[0,\infty)$. Aidon kasvavuuden näkee tekijöihinjaosta:
  kun $0\le u<v$,
  \[
    v^n-u^n=(v-u)\sum_{j=0}^{n-1}v^{n-1-j}u^j>0.
  \]
  Summan ensimmäinen termi on positiivinen (ja kun $n=1$,
  summa on $1$).
  Jos $x=0$, ratkaisu on $u=0$. Jos $x>0$, valitaan $M=x+1>1$.
  Tällöin $P_n(0)=0<x<M^n=P_n(M)$, joten väliarvolause
  antaa $u\in(0,M)$, jolla $u^n=x$.
  Aito kasvavuus takaa ratkaisun yksikäsitteisyyden.
  Siten $P_n$ on bijektio väliltä $[0,\infty)$ itselleen.
  Sen käänteisfunktion jatkuvuus ja aito kasvavuus seuraavat
  \cref{thm:inverse-continuity}.
\end{proof}

\subsubsection*{Rationaaliset potenssit}

\begin{definition}[Rationaalinen potenssi]
  Kun $x>0$, $m\ge0$ on kokonaisluku ja $n\ge1$, määritellään
  \[
    x^{m/n}=\left(x^{1/n}\right)^m.
  \]
  Erityisesti $x^0=1$. Negatiivisille rationaalisille eksponenteille
  asetetaan $x^{-r}=1/x^r$, kun $r>0$.
\end{definition}

\begin{lemma}[Määritelmän riippumattomuus esitystavasta]
  Jos $m/n=p/q$, missä $m,p\ge0$ ja $n,q\ge1$ ovat
  kokonaislukuja, niin
  \[
    \left(x^{1/n}\right)^m=\left(x^{1/q}\right)^p.
  \]
\end{lemma}
\begin{proof}
  Molemmat puolet ovat positiivisia, ja niiden $nq$:nnet
  potenssit ovat vastaavasti $x^{mq}$ ja $x^{pn}$.
  Koska $mq=pn$, nämä potenssit ovat samat.
  Potenssifunktion $u\mapsto u^{nq}$ injektiivisyys
  positiivisilla luvuilla antaa väitteen.
  Negatiivisten eksponenttien tapauksessa otetaan käänteisluvut.
\end{proof}

\newpage
\begin{theorem}[Potenssien laskusäännöt]
  Olkoot $x,y>0$, $m,p\ge0$ ja $n,q\ge1$ kokonaislukuja.
  Tällöin
  \begin{enumerate}[label=\textup{(\alph*)}]
    \item $\displaystyle (x^m)^{1/n}=(x^{1/n})^m=x^{m/n}$,
    \item $\displaystyle x^{m/n+p/q}=x^{m/n}x^{p/q}$,
    \item $\displaystyle (xy)^{1/n}=x^{1/n}y^{1/n}$.
  \end{enumerate}
\end{theorem}
\begin{proof}
  \textbf{(a)} Kokonaislukupotenssien laskusäännöillä
  \[
    \left((x^{1/n})^m\right)^n
    =\left((x^{1/n})^n\right)^m=x^m.
  \]
  Koska $(x^{1/n})^m>0$, juuren yksikäsitteisyys antaa väitteen.

  \smallskip
  \noindent\textbf{(b)} Yhteisen nimittäjän avulla
  $m/n+p/q=(mq+np)/(nq)$, joten
  \[
    \left(x^{m/n+p/q}\right)^{nq}=x^{mq+np}.
  \]
  Toisaalta
  \[
    \left(x^{m/n}x^{p/q}\right)^{nq}
    =\left(x^{m/n}\right)^{nq}\left(x^{p/q}\right)^{nq}
    =x^{mq}x^{np}=x^{mq+np}.
  \]
  Positiivisten lukujen $nq$:nnet potenssit ovat samat,
  joten myös itse luvut ovat samat.

  \smallskip
  \noindent\textbf{(c)} Koska
  \[
    \left(x^{1/n}y^{1/n}\right)^n
    =\left(x^{1/n}\right)^n\left(y^{1/n}\right)^n=xy
  \]
  ja $x^{1/n}y^{1/n}>0$, juuren yksikäsitteisyys antaa väitteen.
\end{proof}

\begin{corollary}[Rationaalisten potenssifunktioiden jatkuvuus]
  Jokaisella $r\in\QQ$ funktio $x\mapsto x^r$ on jatkuva
  välillä $(0,\infty)$.
\end{corollary}
\begin{proof}
  Jos $r=m/n\ge0$, funktio on jatkuvan juurifunktion ja
  jatkuvan kokonaislukupotenssin yhdistetty funktio.
  Jos $r<0$, käytetään esitystä $x^r=1/x^{-r}$ sekä
  osamäärän jatkuvuutta; nimittäjä on positiivinen.
\end{proof}

\begin{remark}
  Kun $r>0$ on rationaalinen, määritellään lisäksi $0^r=0$.
  Sama juurifunktion ja potenssifunktion yhdistäminen todistaa
  jatkuvuuden myös välillä $[0,\infty)$.
  Negatiivista potenssia ei määritellä nollassa.
\end{remark}

\newpage
\subsection{Sinin ja kosinin jatkuvuus}

\begin{theorem}
  Funktiot $\sin\colon\RR\to\RR$ ja
  $\cos\colon\RR\to\RR$ ovat jatkuvia.
\end{theorem}
\begin{proof}
  \textbf{1. Sini origossa.}
  Yksikköympyrän geometriasta tiedetään, että radiaaneina
  mitatulle kulmalle $0<t<\pi/2$ pätee
  \[
    0<\sin t<t.
  \]
  Ylemmän arvion voi nähdä vertaamalla yksikköympyrän
  keskuskulman $t$ määräämää kolmiota sektoriin:
  niiden pinta-alat ovat $\sin t/2$ ja $t/2$.
  Puristusperiaatteen nojalla $\lim_{t\to0^+}\sin t=0$.
  Koska $\sin(-t)=-\sin t$, myös
  \[
    \lim_{x\to0^-}\sin x
    =\lim_{t\to0^+}(-\sin t)=0.
  \]
  Siis $\lim_{x\to0}\sin x=0=\sin0$, eli sini on jatkuva
  pisteessä $0$.

  \medskip
  \noindent\textbf{2. Kosini origossa.}
  Kun $-\pi/2<x<\pi/2$, kosini on positiivinen ja
  \[
    \cos x=\sqrt{1-\sin^2x}.
  \]
  Sinin jatkuvuus pisteessä $0$, jatkuvuuden laskusäännöt
  ja neliöjuurifunktion jatkuvuus antavat
  \[
    \lim_{x\to0}\cos x=\sqrt{1-0^2}=1=\cos0.
  \]
  Siis myös kosini on jatkuva pisteessä $0$.

  \medskip
  \noindent\textbf{3. Sini ja kosini mielivaltaisessa pisteessä.}
  Kiinnitetään $x_0\in\RR$ ja asetetaan $h=x-x_0$.
  Tällöin $x\to x_0$ täsmälleen silloin, kun $h\to0$.
  Yhteenlaskukaavoista ja jo todistetuista raja-arvoista saadaan
  \begin{align*}
    \lim_{x\to x_0}\sin x
    &=\lim_{h\to0}\sin(x_0+h)\\
    &=\lim_{h\to0}(\sin x_0\cos h+\cos x_0\sin h)\\
    &=\sin x_0\cdot1+\cos x_0\cdot0=\sin x_0,
  \end{align*}
  sekä
  \begin{align*}
    \lim_{x\to x_0}\cos x
    &=\lim_{h\to0}\cos(x_0+h)\\
    &=\lim_{h\to0}(\cos x_0\cos h-\sin x_0\sin h)\\
    &=\cos x_0\cdot1-\sin x_0\cdot0=\cos x_0.
  \end{align*}
  Koska $x_0$ oli mielivaltainen, molemmat funktiot ovat
  jatkuvia koko reaalilukujoukossa.
\end{proof}

\end{document}
